\documentclass{article}
\usepackage{oz}

\begin{document}

% ============================================================
% Library Management System - Z Specification
% Author: Anton Zagzin
% Course: Software Engineering Models and Methods (SEMM), Lab 4
% Tool: CZT (Community Z Tools) 1.6
% ============================================================

\begin{zsection}
  \SECTION LibrarySystem \parents standard\_toolkit
\end{zsection}

% ============================================================
% 1. BASIC TYPES AND FREE TYPES
% ============================================================

% Given sets representing the core domain entities
\begin{zed}
  [PERSON, ISBN, TITLE]
\end{zed}

% Genre classification for books
\begin{zed}
  GENRE ::= fiction | nonfiction | science | history | biography
\end{zed}

% Response codes for operations
\begin{zed}
  REPORT ::= ok \\
    | bookNotFound \\
    | memberNotFound \\
    | bookAlreadyExists \\
    | memberAlreadyExists \\
    | alreadyBorrowed \\
    | notBorrowed \\
    | loanLimitReached
\end{zed}

% Maximum number of books a single member can borrow
\begin{axdef}
  maxLoans : \nat_1
\where
  maxLoans = 5
\end{axdef}

% ============================================================
% 2. STATE SCHEMA (6 state variables, 6 invariants)
% ============================================================

% The main state of the library system
% State variables:
%   books    - set of all book ISBNs in the library
%   members  - set of registered library members
%   title    - mapping from ISBN to book title
%   genre    - mapping from ISBN to book genre
%   loans    - current active loans (book -> borrower)
%   waitlist - reservation queue per book
\begin{schema}{Library}
  books : \power ISBN \\
  members : \power PERSON \\
  title : ISBN \pfun TITLE \\
  genre : ISBN \pfun GENRE \\
  loans : ISBN \pfun PERSON \\
  waitlist : ISBN \pfun \seq PERSON
\where
  % Inv1: every book has a title
  \dom title = books \\
  % Inv2: every book has a genre
  \dom genre = books \\
  % Inv3: only registered books can be loaned out
  \dom loans \subseteq books \\
  % Inv4: only registered members can borrow
  \ran loans \subseteq members \\
  % Inv5: waitlists exist only for registered books
  \dom waitlist \subseteq books \\
  % Inv6: no member may borrow more than maxLoans books
  \forall m : members @ \# \{ b : \dom loans | loans(b) = m \} \leq maxLoans
\end{schema}

% ============================================================
% 3. FUNCTIONS (4 axiomatic definitions)
% ============================================================

% Function 1: books currently borrowed by a given member
\begin{axdef}
  borrowedBy : PERSON \cross \power (ISBN \cross PERSON) \fun \power ISBN
\where
  \forall m : PERSON; ln : \power (ISBN \cross PERSON) @ \\
    \t1 borrowedBy(m, ln) = \{ b : \dom ln | (b, m) \in ln \}
\end{axdef}

% Function 2: set of books currently available (not on loan)
\begin{axdef}
  availableBooks : \power ISBN \cross \power (ISBN \cross PERSON) \fun \power ISBN
\where
  \forall bs : \power ISBN; ln : \power (ISBN \cross PERSON) @ \\
    \t1 availableBooks(bs, ln) = bs \setminus \dom ln
\end{axdef}

% Function 3: number of loans for a given member
\begin{axdef}
  loanCount : PERSON \cross \power (ISBN \cross PERSON) \fun \nat
\where
  \forall m : PERSON; ln : \power (ISBN \cross PERSON) @ \\
    \t1 loanCount(m, ln) = \# \{ b : \dom ln | (b, m) \in ln \}
\end{axdef}

% Function 4: filter books by genre
\begin{axdef}
  booksByGenre : GENRE \cross \power ISBN \cross (ISBN \pfun GENRE) \fun \power ISBN
\where
  \forall g : GENRE; bs : \power ISBN; gn : ISBN \pfun GENRE @ \\
    \t1 booksByGenre(g, bs, gn) = \{ b : bs | b \in \dom gn \land gn(b) = g \}
\end{axdef}

% ============================================================
% 4. INITIAL STATE
% ============================================================

% The library starts empty
\begin{schema}{InitLibrary}
  Library~'
\where
  books' = \emptyset \\
  members' = \emptyset \\
  title' = \emptyset \\
  genre' = \emptyset \\
  loans' = \emptyset \\
  waitlist' = \emptyset
\end{schema}

% ============================================================
% 5. OPERATIONS (7 operations, all with pre and post conditions)
% ============================================================

% --- Operation 1: Add a new book to the library ---
\begin{schema}{AddBook}
  \Delta Library \\
  isbn? : ISBN \\
  t? : TITLE \\
  g? : GENRE
\where
  % Pre: the book must not already exist
  isbn? \notin books \\
  % Post: book is added with its title and genre
  books' = books \cup \{ isbn? \} \\
  title' = title \cup \{ isbn? \mapsto t? \} \\
  genre' = genre \cup \{ isbn? \mapsto g? \} \\
  waitlist' = waitlist \cup \{ isbn? \mapsto \langle \rangle \} \\
  % Frame: loans and members unchanged
  loans' = loans \\
  members' = members
\end{schema}

% --- Operation 2: Remove a book from the library ---
\begin{schema}{RemoveBook}
  \Delta Library \\
  isbn? : ISBN
\where
  % Pre: book exists and is not currently on loan
  isbn? \in books \\
  isbn? \notin \dom loans \\
  % Post: book and all its mappings are removed
  books' = books \setminus \{ isbn? \} \\
  title' = \{ isbn? \} \ndres title \\
  genre' = \{ isbn? \} \ndres genre \\
  waitlist' = \{ isbn? \} \ndres waitlist \\
  % Frame
  loans' = loans \\
  members' = members
\end{schema}

% --- Operation 3: Register a new member ---
\begin{schema}{RegisterMember}
  \Delta Library \\
  p? : PERSON
\where
  % Pre: person is not already registered
  p? \notin members \\
  % Post: person is added to the member set
  members' = members \cup \{ p? \} \\
  % Frame
  books' = books \\
  title' = title \\
  genre' = genre \\
  loans' = loans \\
  waitlist' = waitlist
\end{schema}

% --- Operation 4: Remove a member ---
\begin{schema}{RemoveMember}
  \Delta Library \\
  p? : PERSON
\where
  % Pre: member exists and has no active loans
  p? \in members \\
  p? \notin \ran loans \\
  % Post: member is removed
  members' = members \setminus \{ p? \} \\
  % Frame
  books' = books \\
  title' = title \\
  genre' = genre \\
  loans' = loans \\
  waitlist' = waitlist
\end{schema}

% --- Operation 5: Borrow a book ---
\begin{schema}{BorrowBook}
  \Delta Library \\
  isbn? : ISBN \\
  p? : PERSON
\where
  % Pre: book exists, member exists, book is available, member under limit
  isbn? \in books \\
  p? \in members \\
  isbn? \notin \dom loans \\
  loanCount(p?, loans) < maxLoans \\
  % Post: loan is recorded
  loans' = loans \cup \{ isbn? \mapsto p? \} \\
  % Frame
  books' = books \\
  members' = members \\
  title' = title \\
  genre' = genre \\
  waitlist' = waitlist
\end{schema}

% --- Operation 6: Return a book ---
\begin{schema}{ReturnBook}
  \Delta Library \\
  isbn? : ISBN \\
  p? : PERSON
\where
  % Pre: book is currently borrowed by this member
  isbn? \in \dom loans \\
  loans(isbn?) = p? \\
  % Post: loan record is removed
  loans' = \{ isbn? \} \ndres loans \\
  % Frame
  books' = books \\
  members' = members \\
  title' = title \\
  genre' = genre \\
  waitlist' = waitlist
\end{schema}

% --- Operation 7: Search books by genre (returns a set) ---
\begin{schema}{SearchByGenre}
  \Xi Library \\
  g? : GENRE \\
  result! : \power ISBN
\where
  % Pre: library has books
  books \neq \emptyset \\
  % Post: returns all books of the requested genre
  result! = booksByGenre(g?, books, genre)
\end{schema}

% --- Operation 8: Get all books borrowed by a member (returns a set) ---
\begin{schema}{GetMemberBooks}
  \Xi Library \\
  p? : PERSON \\
  result! : \power ISBN
\where
  % Pre: person is a registered member
  p? \in members \\
  % Post: returns the set of books currently borrowed by this member
  result! = borrowedBy(p?, loans)
\end{schema}

% ============================================================
% 6. ERROR SCHEMAS AND SCHEMA CALCULUS
% ============================================================

% --- Success and error report schemas ---
\begin{schema}{Success}
  r! : REPORT
\where
  r! = ok
\end{schema}

\begin{schema}{BookNotFoundError}
  \Xi Library \\
  r! : REPORT
\where
  r! = bookNotFound
\end{schema}

\begin{schema}{MemberNotFoundError}
  \Xi Library \\
  r! : REPORT
\where
  r! = memberNotFound
\end{schema}

\begin{schema}{BookAlreadyExistsError}
  \Xi Library \\
  r! : REPORT \\
  isbn? : ISBN
\where
  isbn? \in books \\
  r! = bookAlreadyExists
\end{schema}

\begin{schema}{MemberAlreadyExistsError}
  \Xi Library \\
  r! : REPORT \\
  p? : PERSON
\where
  p? \in members \\
  r! = memberAlreadyExists
\end{schema}

\begin{schema}{AlreadyBorrowedError}
  \Xi Library \\
  r! : REPORT \\
  isbn? : ISBN
\where
  isbn? \in \dom loans \\
  r! = alreadyBorrowed
\end{schema}

\begin{schema}{NotBorrowedError}
  \Xi Library \\
  r! : REPORT \\
  isbn? : ISBN
\where
  isbn? \notin \dom loans \\
  r! = notBorrowed
\end{schema}

\begin{schema}{LoanLimitError}
  \Xi Library \\
  r! : REPORT \\
  p? : PERSON
\where
  loanCount(p?, loans) \geq maxLoans \\
  r! = loanLimitReached
\end{schema}

% ============================================================
% Schema Calculus Operator 1: CONJUNCTION (and)
% Combine AddBook with Success to form a robust operation
% ============================================================
\begin{zed}
  AddBookOk == AddBook \land Success
\end{zed}

% ============================================================
% Schema Calculus Operator 2: DISJUNCTION (or)
% Total operation: either succeeds or returns an error
% ============================================================
\begin{zed}
  TotalAddBook == AddBookOk \lor BookAlreadyExistsError
\end{zed}

% ============================================================
% Schema Calculus Operator 3: NEGATION (not)
% BookNotOnLoan describes the condition when a book is not loaned
% ============================================================
\begin{schema}{BookOnLoan}
  Library \\
  isbn? : ISBN
\where
  isbn? \in \dom loans
\end{schema}

\begin{zed}
  BookNotOnLoan == \lnot BookOnLoan
\end{zed}

% ============================================================
% Schema Calculus Operator 4: COMPOSITION (sequential)
% Borrow a book and then immediately query the member's books
% ============================================================
\begin{zed}
  BorrowAndQuery == BorrowBook \semi GetMemberBooks
\end{zed}

% Additional total operations using schema calculus (disjunction)
\begin{zed}
  BorrowBookOk == BorrowBook \land Success
\end{zed}

\begin{zed}
  TotalBorrowBook == BorrowBookOk \lor BookNotFoundError \\
    \t1 \lor AlreadyBorrowedError \lor LoanLimitError
\end{zed}

\begin{zed}
  ReturnBookOk == ReturnBook \land Success
\end{zed}

\begin{zed}
  TotalReturnBook == ReturnBookOk \lor NotBorrowedError
\end{zed}

% ============================================================
% 7. DATA REFINEMENT
% Abstract (set-based) -> Concrete (sequence-based)
% ============================================================

% --- Concrete state: books and loans stored as sequences ---
\begin{schema}{ConcreteLibrary}
  cbooks : \seq ISBN \\
  cmembers : \seq PERSON \\
  ctitle : \seq (ISBN \cross TITLE) \\
  cgenre : \seq (ISBN \cross GENRE) \\
  cloans : \seq (ISBN \cross PERSON) \\
  cwaitlist : \seq (ISBN \cross \seq PERSON)
\where
  % Concrete invariants mirror abstract ones over sequences
  % All ISBNs in cbooks are unique
  \forall i, j : \dom cbooks @ \\
    \t1 i \neq j \implies cbooks(i) \neq cbooks(j) \\
  % All members are unique
  \forall i, j : \dom cmembers @ \\
    \t1 i \neq j \implies cmembers(i) \neq cmembers(j) \\
  % Loaned books must be in the book collection
  \forall i : \dom cloans @ \\
    \t1 first(cloans(i)) \in \ran cbooks \\
  % Loan borrowers must be registered members
  \forall i : \dom cloans @ \\
    \t1 second(cloans(i)) \in \ran cmembers \\
  % Loan limit applies
  \forall m : \ran cmembers @ \\
    \t1 \# \{ i : \dom cloans | second(cloans(i)) = m \} \leq maxLoans
\end{schema}

% --- Retrieve relation: maps concrete state to abstract state ---
\begin{schema}{Retrieve}
  Library \\
  ConcreteLibrary
\where
  books = \ran cbooks \\
  members = \ran cmembers \\
  title = \ran ctitle \\
  genre = \ran cgenre \\
  loans = \ran cloans \\
  waitlist = \ran cwaitlist
\end{schema}

% --- Concrete initial state ---
\begin{schema}{ConcreteInitLibrary}
  ConcreteLibrary~'
\where
  cbooks' = \langle \rangle \\
  cmembers' = \langle \rangle \\
  ctitle' = \langle \rangle \\
  cgenre' = \langle \rangle \\
  cloans' = \langle \rangle \\
  cwaitlist' = \langle \rangle
\end{schema}

% --- Refined Operation 1: Concrete BorrowBook ---
% Uses sequence append instead of set union
\begin{schema}{ConcreteBorrowBook}
  \Delta ConcreteLibrary \\
  isbn? : ISBN \\
  p? : PERSON
\where
  % Pre: book exists in sequence, member exists, book not already loaned
  isbn? \in \ran cbooks \\
  p? \in \ran cmembers \\
  \lnot (\exists i : \dom cloans @ first(cloans(i)) = isbn?) \\
  \# \{ i : \dom cloans | second(cloans(i)) = p? \} < maxLoans \\
  % Post: append new loan record to the sequence
  cloans' = cloans \cat \langle (isbn?, p?) \rangle \\
  % Frame
  cbooks' = cbooks \\
  cmembers' = cmembers \\
  ctitle' = ctitle \\
  cgenre' = cgenre \\
  cwaitlist' = cwaitlist
\end{schema}

% --- Refined Operation 2: Concrete ReturnBook ---
% Filters the loan sequence to remove the returned book
\begin{schema}{ConcreteReturnBook}
  \Delta ConcreteLibrary \\
  isbn? : ISBN \\
  p? : PERSON
\where
  % Pre: the loan record exists
  \exists i : \dom cloans @ cloans(i) = (isbn?, p?) \\
  % Post: remove the loan entry by filtering
  \ran cloans' = \ran cloans \setminus \{ (isbn?, p?) \} \\
  % All other entries preserved
  \forall i : \dom cloans' @ cloans'(i) \in \ran cloans \setminus \{ (isbn?, p?) \} \\
  % Frame
  cbooks' = cbooks \\
  cmembers' = cmembers \\
  ctitle' = ctitle \\
  cgenre' = cgenre \\
  cwaitlist' = cwaitlist
\end{schema}

% --- Refined Operation 3: Concrete SearchByGenre ---
% Iterates over genre sequence instead of using set comprehension
\begin{schema}{ConcreteSearchByGenre}
  \Xi ConcreteLibrary \\
  g? : GENRE \\
  result! : \power ISBN
\where
  % Pre: there are books
  \ran cbooks \neq \emptyset \\
  % Post: collect ISBNs whose genre matches
  result! = \{ i : \dom cgenre | second(cgenre(i)) = g? @ first(cgenre(i)) \}
\end{schema}

\end{document}
